%============================================================
% CH1 — CHAPTER SUMMARY & WHAT COMES NEXT   (tag g1s)
%   cheat sheet, the five patterns in pictures, storage picture,
%   "which pattern?" recognition table, complexity tables, mind map,
%   pattern growth map, self-check quiz (answers at the end of the chapter),
%   readiness checklist, final words
%============================================================

\subsection{The Chapter on One Page}
\label{ssec:ch1-summary-main}

\begin{chaptertheorybox}{Cheat Sheet --- Chapter 1}
\begin{center}
\small
\renewcommand{\arraystretch}{1.45}
\begin{tabularx}{\linewidth}{@{}>{\bfseries}l >{\raggedright\arraybackslash}p{3.3cm} >{\raggedright\arraybackslash}X l@{}}
\toprule
Pattern & Graph model & Core move & Cost \\
\midrule
P1 Degrees & edge list, one counter per vertex & \texttt{deg[u]++}, \texttt{deg[v]++} (or in/out for arrows), then a formula on the counters & \bigO{n+m} \\
P2 Reachability & adjacency list & DFS or BFS from the source; mark a vertex \emph{when it is discovered} & \bigO{n+m} \\
P3 Grids & cell $(r,c)$, 4 neighbours & flood fill from every free, unvisited cell; one fill $=$ one region & \bigO{RC} \\
P4 Implicit states & state $=$ number / tuple, edge $=$ one move & BFS over states, neighbours generated on the fly & \bigO{S \cdot b} \\
P5 Functional & \texttt{nxt[v]}, out-degree $1$ & follow arrows; colours $0/1/2$ (or timestamps) reveal the cycle; Floyd's two pointers when memory is \bigO{1} & \bigO{n} \\
\bottomrule
\end{tabularx}
\end{center}

\textbf{The vocabulary in one breath:} $\sum_v \deg(v) = 2|E|$, so the number of odd-degree vertices is even; a tree is connected with $n-1$ edges; a DAG has no directed cycle; every component of a functional graph is one cycle with trees hanging into it.

\textbf{The storage rule:} adjacency list by default; an adjacency matrix only when $n \lesssim 5000$; an undirected edge is stored as two directed edges; reset everything between test cases.

\textbf{The visited rule:} every vertex enters the stack or queue at most once, every edge is looked at a constant number of times $\Rightarrow$ \bigO{|V|+|E|}.

\textbf{The golden rule:} before writing code, fill in the Graph Model --- what is a vertex, what is an edge, directed or not, weighted or not, what is the question in graph language, and how large are $|V|$ and $|E|$?
\end{chaptertheorybox}

\subsection{The Five Patterns in Pictures}

Each panel shows the smallest picture that still explains its pattern. If you can redraw these six panels from memory, you remember the chapter.

\begin{diagrambox}
\begin{tikzpicture}[
    sv/.style={minimum size=6mm, font=\scriptsize\bfseries},
    pnl/.style={draw=patternborder, fill=patternbg!45, rounded corners=3pt, line width=0.8pt},
    ptitle/.style={font=\scriptsize\bfseries\color{patternborder!80!black}},
    pcap/.style={font=\scriptsize, text width=4.5cm, align=center, text=darkgray},
    mv/.style={font=\tiny\bfseries, fill=white, inner sep=0.8pt, text=deepblue}]
    % ---------------- P1 degrees ----------------
    \begin{scope}[shift={(0,4.6)}]
        \draw[pnl] (0,0) rectangle (4.8,4.2);
        \node[ptitle] at (2.4,3.9) {P1 Degrees \& Counting};
        \node[gv, sv] (a) at (0.8,3.15) {1};
        \node[gv, sv] (b) at (4.0,3.15) {2};
        \node[gv, sv] (c) at (0.8,1.65) {4};
        \node[gvdst, sv] (j) at (2.4,2.3) {3};
        \draw[gde] (a) -- (j); \draw[gde] (b) -- (j); \draw[gde] (c) -- (j); \draw[gde] (a) -- (b);
        \node[gdist, anchor=west, text=accentorange] at (2.8,2.0) {in $3$, out $0$};
        \node[pcap] at (2.4,0.55) {One pass counts in/out-degrees;\\judge: in $n{-}1$, out $0$.};
    \end{scope}
    % ---------------- P2 reachability ----------------
    \begin{scope}[shift={(5.1,4.6)}]
        \draw[pnl] (0,0) rectangle (4.8,4.2);
        \node[ptitle] at (2.4,3.9) {P2 Reachability};
        \node[gvsrc, sv] (s) at (0.7,2.9) {0};
        \node[gvdone, sv] (u1) at (1.9,3.15) {1};
        \node[gvdone, sv] (u2) at (1.2,1.7) {2};
        \node[gvdone, sv] (u3) at (2.6,2.1) {3};
        \node[gvdead, sv] (x1) at (3.9,3.0) {4};
        \node[gvdead, sv] (x2) at (4.0,1.7) {5};
        \draw[gtree] (s) -- (u1); \draw[gtree] (s) -- (u2); \draw[gtree] (u1) -- (u3);
        \draw[ge] (u2) -- (u3); \draw[gfaded] (x1) -- (x2);
        \node[pcap] at (2.4,0.55) {DFS/BFS from $0$ with \texttt{visited}:\\$\{0,1,2,3\}$ reached, $\{4,5\}$ not.};
    \end{scope}
    % ---------------- P3 grids ----------------
    \begin{scope}[shift={(10.2,4.6)}]
        \draw[pnl] (0,0) rectangle (4.8,4.2);
        \node[ptitle] at (2.4,3.9) {P3 Grids as Graphs};
        % cell (column c, row r) has its lower-left corner at (0.75 + 0.55c, 2.75 - 0.55r)
        \foreach \c/\r in {0/0,3/0,0/1,2/1,3/1,5/1,2/2,5/2,0/3,1/3,2/3,4/3,5/3}
            \fill[blockedfill] (0.75+0.55*\c, 2.75-0.55*\r) rectangle ++(0.55,0.55);
        \foreach \c/\r in {1/0,2/0,1/1,0/2,1/2}
            \fill[primaryblue!22] (0.75+0.55*\c, 2.75-0.55*\r) rectangle ++(0.55,0.55);
        \foreach \c/\r in {4/0,5/0,4/1,3/2,4/2,3/3}
            \fill[dupfill] (0.75+0.55*\c, 2.75-0.55*\r) rectangle ++(0.55,0.55);
        \foreach \k in {0,...,6} \draw[darkgray!50] (0.75+0.55*\k, 1.1) -- (0.75+0.55*\k, 3.3);
        \foreach \k in {0,...,4} \draw[darkgray!50] (0.75, 1.1+0.55*\k) -- (4.05, 1.1+0.55*\k);
        \node[font=\scriptsize\bfseries, text=primaryblue] at (0.75+0.55*1.5, 2.75-0.55*0.5) {A};
        \node[font=\scriptsize\bfseries, text=recursionred] at (0.75+0.55*4.5, 2.75-0.55*0.5) {B};
        \node[pcap] at (2.4,0.55) {Cells are vertices, sides are edges;\\one flood fill per room: $2$ rooms.};
    \end{scope}
    % ---------------- P4 implicit states ----------------
    \begin{scope}[shift={(0,0)}]
        \draw[pnl] (0,0) rectangle (4.8,4.2);
        \node[ptitle] at (2.4,3.9) {P4 Implicit State Graphs};
        \node[gvsrc, sv] (n4) at (2.4,3.25) {4};
        \node[gvseen, sv] (n8) at (1.3,2.4) {8};
        \node[gvseen, sv] (n3) at (3.5,2.4) {3};
        \node[gv, sv] (n16) at (0.6,1.55) {16};
        \node[gv, sv] (n7) at (1.8,1.55) {7};
        \node[gvdst, sv] (n6) at (3.0,1.55) {6};
        \node[gv, sv] (n2) at (4.2,1.55) {2};
        \draw[gde] (n4) -- node[mv, pos=0.45] {$\times2$} (n8);
        \draw[gpath, -{Stealth[length=2.4mm]}] (n4) -- node[mv, pos=0.45] {$-1$} (n3);
        \draw[gde] (n8) -- (n16); \draw[gde] (n8) -- (n7);
        \draw[gpath, -{Stealth[length=2.4mm]}] (n3) -- node[mv, pos=0.45] {$\times2$} (n6);
        \draw[gde] (n3) -- (n2);
        \node[pcap] at (2.4,0.55) {Numbers $=$ vertices, presses $=$ edges;\\BFS: $4 \to 3 \to 6$, $2$ presses.};
    \end{scope}
    % ---------------- P5 functional ----------------
    \begin{scope}[shift={(5.1,0)}]
        \draw[pnl] (0,0) rectangle (4.8,4.2);
        \node[ptitle] at (2.4,3.9) {P5 Functional Graphs};
        \node[gv, sv] (f1) at (0.6,3.05) {1};
        \node[gv, sv] (f6) at (0.6,1.65) {6};
        \node[gv, sv] (f2) at (1.55,2.25) {2};
        \node[gvsel, sv] (f3) at (2.65,2.25) {3};
        \node[gvsel, sv] (f4) at (3.85,3.0) {4};
        \node[gvsel, sv] (f5) at (3.85,1.6) {5};
        \draw[gde] (f1) -- (f2); \draw[gde] (f6) -- (f2); \draw[gde] (f2) -- (f3);
        \draw[gde, softpurple] (f3) -- (f4); \draw[gde, softpurple] (f4) -- (f5); \draw[gde, softpurple] (f5) -- (f3);
        \node[pcap] at (2.4,0.55) {One arrow leaves every vertex;\\every walk ends in the cycle $3,4,5$.};
    \end{scope}
    % ---------------- the Graph Model ----------------
    \begin{scope}[shift={(10.2,0)}]
        \draw[pnl, draw=definitionborder, fill=definitionbg] (0,0) rectangle (4.8,4.2);
        \node[ptitle, text=definitionborder] at (2.4,3.9) {Before Any Code: the Graph Model};
        \node[font=\scriptsize, text width=4.3cm, align=left, anchor=north west] at (0.25,3.6) {%
            \textbf{1.} What is a vertex?\\
            \textbf{2.} What is an edge? Directed?\\
            \textbf{3.} Weighted? Parallel edges?\\
            \textbf{4.} The question in graph words\\
            \textbf{5.} $|V|$ and $|E|$ from the limits\\
            \textbf{6.} Shape: tree, DAG, grid,\\
            \phantom{\textbf{6.} }out-degree $1$?};
        \node[pcap] at (2.4,0.55) {Fill it in first;\\the pattern usually follows.};
    \end{scope}
\end{tikzpicture}
\figcap{The five patterns of the chapter and the one habit that connects them. Colours: green border $=$ source, orange $=$ target, blue fill $=$ finished, grey $=$ unreachable, purple $=$ on the cycle, thick orange $=$ chosen path.}
\end{diagrambox}

\subsection{Three Ways to Store the Same Graph}

\begin{diagrambox}
\begin{tikzpicture}[
    sv/.style={minimum size=6.5mm, font=\scriptsize\bfseries},
    mc/.style={draw=darkgray!60, minimum size=5.5mm, inner sep=0pt, font=\scriptsize},
    mh/.style={minimum size=5.5mm, inner sep=0pt, font=\scriptsize\bfseries\color{deepblue}},
    hd/.style={font=\scriptsize\bfseries\color{deepblue}}]
    % the graph
    \node[hd] at (1.1,2.75) {the graph};
    \node[gv, sv] (v0) at (0.3,1.9) {0};
    \node[gv, sv] (v1) at (1.9,1.9) {1};
    \node[gv, sv] (v2) at (1.1,0.6) {2};
    \node[gv, sv] (v3) at (2.4,-0.3) {3};
    \draw[ge] (v0) -- (v1); \draw[ge] (v0) -- (v2); \draw[ge] (v1) -- (v2); \draw[ge] (v2) -- (v3);
    \node[font=\tiny, text=darkgray] at (1.1,-0.85) {degrees $2,2,3,1$; sum $8 = 2 \cdot 4$};
    % edge list
    \node[hd] at (4.1,2.75) {edge list};
    \foreach \e [count=\k from 0] in {{(0,1)},{(0,2)},{(1,2)},{(2,3)}}
        \node[mc, minimum width=12mm] at (4.1,1.9-0.6*\k) {\e};
    \node[font=\tiny, text=darkgray] at (4.1,-0.85) {$m$ pairs};
    % adjacency matrix
    \node[hd] at (7.25,2.75) {adjacency matrix};
    \foreach \j in {0,...,3} \node[mh] at (6.35+0.55*\j,2.25) {\j};
    \foreach \i in {0,...,3} \node[mh] at (5.75,1.7-0.55*\i) {\i};
    \foreach \row [count=\i from 0] in {{0,1,1,0},{1,0,1,0},{1,1,0,1},{0,0,1,0}}
        \foreach \x [count=\j from 0] in \row {
            \ifnum\x=1 \node[mc, fill=primaryblue!18, font=\scriptsize\bfseries] at (6.35+0.55*\j,1.7-0.55*\i) {1};
            \else \node[mc, text=midgray] at (6.35+0.55*\j,1.7-0.55*\i) {0}; \fi}
    \node[font=\tiny, text=darkgray] at (7.2,-0.85) {$n \times n$ cells, symmetric};
    % adjacency list
    \node[hd] at (11.6,2.75) {adjacency list};
    \foreach \v/\lst [count=\i from 0] in {0/{1,2}, 1/{0,2}, 2/{0,1,3}, 3/{2}} {
        \node[mh] at (10.1,1.9-0.6*\i) {\v};
        \foreach \w [count=\k from 0] in \lst
            \node[mc, fill=computedfill] at (10.75+0.6*\k,1.9-0.6*\i) {\w};
        \draw[flow] (10.35,1.9-0.6*\i) -- (10.47,1.9-0.6*\i);
    }
    \node[font=\tiny, text=darkgray] at (11.6,-0.85) {$2m = 8$ entries in total};
\end{tikzpicture}
\figcap{One undirected graph, three representations (Lesson~\ref{ssec:ch1-storing}). The adjacency list stores every undirected edge twice, which is the handshake lemma in memory form.}
\end{diagrambox}

\subsection{Which Pattern? A Recognition Table}

Read the statement, find the row whose left column sounds like it, and start from the model in the middle column. The rows are phrased the way contest statements are phrased, not the way textbooks are.

\begin{center}
\small
\renewcommand{\arraystretch}{1.35}
\begin{tabularx}{\linewidth}{@{}>{\raggedright\arraybackslash}p{4.7cm} >{\raggedright\arraybackslash}X c >{\raggedright\arraybackslash}p{2.6cm}@{}}
\toprule
\textbf{The statement says \dots} & \textbf{The graph in disguise} & \textbf{Pattern} & \textbf{Seen in} \\
\midrule
``everybody trusts / beats / points to \dots'', ``find the special person'' & directed graph; the special vertex has a telltale in/out-degree & P1 & Town Judge \\
``rank / importance / score of a pair of cities'' & a formula in $\deg(u)$, $\deg(v)$ and ``is $u$--$v$ an edge?'' & P1 & Network Rank \\
``$n$ vertices and $n-1$ edges'', ``is it a star / a path?'' & a tree; shape tests are degree tests plus at most one traversal & P1 (+P2) & P1.03, P1.08 \\
``is there a route from $A$ to $B$?'', ``can every room be visited?'' & one DFS/BFS from the source, then read \texttt{visited} & P2 & Path Exists, Keys and Rooms \\
``to do $X$ you first need $Y$'', ``which items are needed?'' & reachability along prerequisite arrows & P2 & P1.12 \\
``map of \texttt{.} and \texttt{\#}'', ``islands'', ``rooms'', ``regions'' & cells $+$ 4 neighbours; count the flood fills & P3 & Counting Rooms \\
``largest / total / border of a region'' & a flood fill that also accumulates something & P3 & Max Area of Island \\
``minimum number of operations / presses to turn $x$ into $y$'' & values are vertices, operations are edges; BFS & P4 & Two Buttons \\
``lock'', ``puzzle'', ``configuration'', ``forbidden states'' & tuples encoded as integers; BFS with a banned set & P4 & Open the Lock \\
``each person names exactly one other'', ``$p_i$'', ``teleport $K$ times'' & functional graph: follow the arrows until a repeat & P5 & Badge \\
``at most one outgoing edge'', ``longest cycle'' & functional graph with dead ends; timestamps per walk & P5 & Longest Cycle \\
``$n+1$ values in $[1,n]$'', ``do not modify the array, \bigO{1} extra space'' & functional graph $i \to a_i$ from vertex $0$; Floyd's tortoise and hare finds the entry point & P5 & Find the Duplicate \\
``a permutation applied again and again'' & disjoint cycles (a functional graph without tails) & P5 & P1.21 \\
\bottomrule
\end{tabularx}
\end{center}

\begin{warningbox}[When the table points to the wrong row]
The words in the statement are only a hint; the constraints decide. ``Minimum number of moves'' with $10^{18}$ possible states is not a BFS (look for a formula or a greedy, as in Two Buttons read backwards), and ``islands'' on a $10^5 \times 10^5$ map are not a flood fill. Always finish the Graph Model, including $|V|$ and $|E|$, before trusting a row.
\end{warningbox}

\subsection{Complexity at a Glance}

\begin{center}
\small
\renewcommand{\arraystretch}{1.35}
\begin{tabularx}{\linewidth}{@{}>{\raggedright\arraybackslash}X c c c c@{}}
\toprule
\textbf{Representation} & \textbf{Memory} & \textbf{Neighbours of $v$} & \textbf{Is $(u,v)$ an edge?} & \textbf{OK for $n = 2\cdot 10^5$?} \\
\midrule
Edge list & \bigO{m} & \bigO{m} & \bigO{m} & \yes\ (for P1 counting) \\
Adjacency matrix & \bigO{n^2} & \bigO{n} & \bigO{1} & \no\ ($4\cdot 10^{10}$ cells) \\
Adjacency list & \bigO{n+m} & \bigO{\deg v} & \bigO{\deg v} & \yes \\
\bottomrule
\end{tabularx}

\medskip
\begin{tabularx}{\linewidth}{@{}l >{\raggedright\arraybackslash}X c c@{}}
\toprule
\textbf{Pattern} & \textbf{Algorithm} & \textbf{Time} & \textbf{Memory} \\
\midrule
P1 Degrees & one pass over the edges, then one pass over the vertices & \bigO{n+m} & \bigO{n} \\
P2 Reachability & DFS or BFS with a \texttt{visited} array & \bigO{n+m} & \bigO{n+m} \\
P3 Grids & flood fill over an $R \times C$ grid ($< 2RC$ edges) & \bigO{RC} & \bigO{RC} \\
P4 Implicit states & BFS over $S$ states with $b$ moves each & \bigO{S \cdot b} & \bigO{S} \\
P5 Functional & follow \texttt{nxt}, each vertex finished once & \bigO{n} & \bigO{n} \\
P5, one start & Floyd's tortoise and hare (Find the Duplicate) & \bigO{n} & \bigO{1} \\
\bottomrule
\end{tabularx}
\end{center}

\noindent Rule of thumb: about $10^8$ simple operations per second. For P4, estimate $S$ \emph{before} coding: $10^4$ lock states are nothing, $10^7$ states are fine with an \texttt{int} array, $10^{12}$ states mean BFS is the wrong tool.

\subsection{Mind Map}

\begin{diagrambox}
\begin{tikzpicture}[
    mindmap, grow cyclic, every node/.style={concept, execute at begin node=\hskip0pt},
    concept color=primaryblue!70, text=white, font=\scriptsize\bfseries,
    level 1/.append style={level distance=4.1cm, sibling angle=72, font=\scriptsize\bfseries},
    level 2/.append style={level distance=2.3cm, sibling angle=40, font=\tiny}]
\node[font=\small\bfseries, minimum size=2.6cm] {Graphs}
    child[concept color=successgreen!80!black] { node {Vocabulary}
        child { node {degree, handshake} }
        child { node {path, cycle, component} }
        child { node {tree, DAG, bipartite} } }
    child[concept color=accentorange] { node {Storage}
        child { node {edge list} }
        child { node {matrix} }
        child { node {adjacency list} } }
    child[concept color=softpurple] { node {Traversal}
        child { node {DFS (stack)} }
        child { node {BFS (queue)} }
        child { node {visited $\Rightarrow$ $\bigO{n{+}m}$} } }
    child[concept color=patternborder] { node {5 Patterns}
        child { node {degrees, reach} }
        child { node {grids, states} }
        child { node {functional} } }
    child[concept color=warningred!80!black] { node {Modelling}
        child { node {what is a vertex?} }
        child { node {what is an edge?} }
        child { node {$|V|$, $|E|$ from limits} } };
\end{tikzpicture}
\end{diagrambox}

\subsection{How Each Pattern Grows Through the Book}

Nothing in this chapter is thrown away. Every later technique is one of these five ideas with a weight, a priority queue, a cleverer order or a richer state.

\begin{diagrambox}
\begin{tikzpicture}[x=3.15cm, y=1.25cm,
    pat/.style={rectangle, rounded corners=3pt, draw=patternborder, fill=patternbg, line width=1pt,
                font=\scriptsize\bfseries, align=center, text width=2.2cm, minimum height=9mm},
    lvl/.style={rectangle, rounded corners=2pt, draw=#1, fill=#1!10, font=\tiny, align=center,
                text width=2.45cm, minimum height=8mm}]
    \node[font=\scriptsize\bfseries, color=deepblue] at (0,0.8) {Chapter 1};
    \node[font=\scriptsize\bfseries, color=successgreen] at (1,0.8) {Part I (Ch.~2--6)};
    \node[font=\scriptsize\bfseries, color=softpurple] at (2,0.8) {Part II--III (Ch.~7--17)};
    \node[font=\scriptsize\bfseries, color=warningred] at (3,0.8) {Part IV--VI (Ch.~18--35)};
    \node[pat] (p1) at (0,0)  {P1 Degrees};
    \node[pat] (p2) at (0,-1) {P2 Reachability};
    \node[pat] (p3) at (0,-2) {P3 Grids};
    \node[pat] (p4) at (0,-3) {P4 State graphs};
    \node[pat] (p5) at (0,-4) {P5 Functional};
    \node[lvl=successgreen] (a1) at (1,0)  {in-degrees and Kahn's order (Ch.~5)};
    \node[lvl=successgreen] (a2) at (1,-1) {DFS, components, DSU (Ch.~2, 4, 6)};
    \node[lvl=successgreen] (a3) at (1,-2) {BFS distances on grids (Ch.~3)};
    \node[lvl=successgreen] (a4) at (1,-3) {BFS on states, multi-source BFS (Ch.~3)};
    \node[lvl=successgreen] (a5) at (1,-4) {cycle detection in DFS (Ch.~2)};
    \node[lvl=softpurple] (b1) at (2,0)  {leaves and degrees in trees (Ch.~13)};
    \node[lvl=softpurple] (b2) at (2,-1) {weighted reachability: Dijkstra (Ch.~7)};
    \node[lvl=softpurple] (b3) at (2,-2) {0-1 BFS on grids (Ch.~8)};
    \node[lvl=softpurple] (b4) at (2,-3) {Dijkstra on states, graph construction (Ch.~7, 12)};
    \node[lvl=softpurple] (b5) at (2,-4) {binary lifting, $K$-th successor (Ch.~11)};
    \node[lvl=warningred] (c1) at (3,0)  {Euler parity, degree tricks (Ch.~21, 22)};
    \node[lvl=warningred] (c2) at (3,-1) {strongly connected components (Ch.~19)};
    \node[lvl=warningred] (c3) at (3,-2) {min cut on grids (Ch.~24)};
    \node[lvl=warningred] (c4) at (3,-3) {flow modelling of states (Ch.~24)};
    \node[lvl=warningred] (c5) at (3,-4) {condensation of SCCs (Ch.~19)};
    \foreach \i in {1,...,5}{\draw[flow] (p\i) -- (a\i); \draw[flow] (a\i) -- (b\i); \draw[flow] (b\i) -- (c\i);}
\end{tikzpicture}
\figcap{Roughly: the further right, the higher the rating. Each arrow means ``the same idea with a richer graph or a cleverer order of visiting it''.}
\end{diagrambox}

\subsection{Self-Check Quiz}

Answer without looking back.
\begin{enumerate}[leftmargin=*]
    \item Can a simple undirected graph have the degree sequence $3, 3, 3, 1, 1$? Why or why not?
    \item An undirected graph with $m$ edges is stored as adjacency lists. How many entries do the lists hold in total?
    \item Why is an adjacency matrix out of the question for $n = 2 \cdot 10^5$, and what do you use instead?
    \item Why does DFS or BFS with a \texttt{visited} array run in \bigO{n+m} and not, say, \bigO{n \cdot m}?
    \item In BFS, why is a vertex marked visited when it is \emph{pushed} into the queue rather than when it is popped?
    \item A recursive DFS crashes on a path of $10^6$ vertices although the code is correct. What happened, and what is the fix?
    \item A connected graph has $n$ vertices and $n-1$ edges. What is it? And with exactly $n$ edges?
    \item In Find the Town Judge, why does counting in- and out-degrees suffice --- no traversal at all?
    \item An $R \times C$ grid with 4-neighbour moves: how many vertices and how many edges does its graph have when no cell is blocked?
    \item Open the Lock: how many states are there, and how many neighbours does each state have?
    \item In Two Buttons ($x \to 2x$ or $x \to x-1$, from $n$ to $m$), why does a BFS never need values above $\max(n, 2m)$?
    \item In a functional graph, what does every connected component look like, and why must every walk eventually repeat a vertex?
\end{enumerate}

\noindent\textit{Answers are at the end of the chapter, after the solution sketches (page~\pageref{ssec:ch1-quiz-answers}). Try every question first.}

\subsection{Ready for Chapter 2?}

\begin{checklistbox}[Move on when every box is ticked]
\begin{itemize}[leftmargin=*]
    \item[$\square$] I can define degree, path, cycle, component, tree, DAG, bipartite and functional graph, and draw a small example of each.
    \item[$\square$] I can read an edge list (0- or 1-indexed, directed or undirected, several test cases) into an adjacency list without looking anything up.
    \item[$\square$] I can write recursive DFS, iterative DFS and BFS with a \texttt{visited} array from memory, each in under 5 minutes.
    \item[$\square$] For every problem of this chapter I can fill in the Graph Model: vertices, edges, direction, the question, $|V|$ and $|E|$.
    \item[$\square$] I can turn a grid, a number game or a lock into a graph and bound its number of states before coding.
    \item[$\square$] I can find the cycle that a walk in a functional graph falls into, and the length of that cycle.
    \item[$\square$] I have solved at least 18 of the 26 practice problems of Tiers 1--3, including at least 3 from Tier 3.
\end{itemize}
\end{checklistbox}

\subsection{What Comes Next}

A graph problem is solved in two steps: \emph{see} the graph, then \emph{walk} it. This chapter trained the first step --- vertices and edges hide in friendships, grids, numbers, locks and arrays of pointers --- and gave you the simplest walks: counting degrees, flooding everything reachable, and following single arrows. Chapter~2 (DFS) takes the walk seriously: entry and exit times, the DFS tree, back edges and cycle detection, and how to avoid the recursion-depth trap for good. Chapter~3 (BFS) turns ``can I get there?'' into ``how fast can I get there?'', with distances, multi-source searches and shortest paths on grids and state graphs. Chapter~4 (components, bipartite), Chapter~5 (DAGs) and Chapter~6 (DSU) then complete the foundations on which every later part of the book is built.
